\section{Finite-description execution and separate resource charges}\label{app:compute}
The Borel program theorem and its finite-bit implementations are distinct assertions. Arbitrary Borel kernels or noncomputable measures do not by themselves provide executable numerical primitives. For a finite implementation assume supplied, parameter-uniform certificates for: rational descriptions of the selected codebook and action probabilities; evaluation of the raw reports and accept regions with controlled probability of an incorrect decision; sampling coins to prescribed variation accuracy; and terminal readout error. A packing or cover used by the geometric theorem must have a computable finite certificate. Products of intervals and Cantor cylinders have such explicit lists. No minimization oracle is presumed to run for free.

\begin{proposition}[Hazard-margin finite execution]\label{prop:finite}
Consider a one-live-label program whose conditional stop probability at every history is at least $p>0$, whose per-call contribution to a loss in $[0,L]$ is at most $b$, and whose mean cost is controlled by its fixed per-call action-cost mean $c$. Suppose a finite-description implementation changes the conditional joint action, retained decision and audit-relevant state law by at most $\eta<p$ in total variation, with the raw report coupled identically, uniformly in every history and model. Then
\[
 \E\tau_{\rm num}\le\frac1{p-\eta},\qquad
 R_{\rm num}\le\frac{b+L\eta}{p-\eta},\qquad
 \E C_{\rm num}\le\frac{c+c_{\max}\eta}{p-\eta}.
\]
Terminal decoder root error adds as in Proposition~\ref{prop:composition}. The same label alphabet suffices. Under the above computational certificates, the geometric matching upper can be implemented by finite code and finite current-call precision without changing its $N,m$ order, while meeting the original mean budget exactly as an inequality.
\end{proposition}
\begin{proof}
This is a decision/routing coupling; it does not require total-variation convergence of a discrete numerical report to the continuous raw report. The implemented stop hazard is at least $p-\eta$, so its actual survival sum is at most $(p-\eta)^{-1}$. Its conditional stopped-loss contribution is at most $b+L\eta$. Summing these two bounds proves the first two assertions; the same argument for expected per-call cost proves the third. These are direct hazard estimates, not expected-time conclusions from a terminal total-variation estimate.

For unit-cost geometric programs at large $N$, build the reference upper with budget $N/2$. Its certified hazard is at least $2/N$. Choose the one-call decision error at most $1/N$ and, more stringently, at most a constant multiple of $F_{d,\alpha}(N,m)/N$. The displayed inequalities then give mean at most $N$ and risk of order $F_{d,\alpha}(N,m)$, since this profile is stable under fixed budget rescaling. Choose readout error of order $\sqrt{F_{d,\alpha}(N,m)}$. Bounded remaining budgets use the fixed one-call decoder. Costs other than one are treated by a strict cost margin in the displayed third inequality. At singular thresholds where rescaling changes the phase, this argument is not used: a specific certified feasible threshold or the exact revealed-atom policy is required instead.
\end{proof}

A program with $m$ centers of coordinate precision $p_r$, finitely many accept-region descriptions of total length $L_A$, and rational action/coin data of length $L_P$ has description length bounded by
$O(kmp_r+L_A+L_P+\log m)$ for a fixed instruction set. Persistent labels require $\lceil\log_2(m+1)\rceil$ bits, not the entire preceding description length. Current-call evaluation requires separately declared workspace $W(p)$ and time $T(p)$. At fixed implemented precision the expected total evaluation time is at most $\E\tau_{\rm num}\,T(p)$, plus preparation, planning and audit costs. A finite expected call count gives no deterministic running-time bound.

For Gaussian tails, overflow regions and threshold-boundary probabilities must both be included in the uniform decision certificate; clipping away their mass is not permitted. The rational tail certificates accompanying Section~\ref{sec:gaussian} certify feasible thresholds and risks, not a cost-free infinite-precision normal report. At a reveal-only zero-risk point, reading the finite report symbol is exact and its finite description is immediate. For a continuous threshold, rational feasible approximants approach the stated value; an exact mathematical threshold is not described as a finite binary real. Program bits, temporary arithmetic space, calibration precision, actual calls, wall-clock primitive time and simulator state therefore remain distinct entries in the resource ledger.
